\subsection{积分曲线}

本节中，X 表示一个 \(C^{r}\) 类流形，\(\xi\) 是 X 上的 \(C^{s}\) 类向量场，其中 r、s 属于 \(N_{K}\)，且 \(s \leqslant r - 1\)。

9.1.1. 设 I 是 K 的开子集，且 \(f: \mathbf{I} \to \mathbf{X}\) 是从 I 到 X 的 \(\mathbf{C}^k\) 类映射（\(k \in \mathbf{N}_{\mathbb{K}}\)，\(k \leqslant r\)）。对于 \(t \in \mathbf{I}\)，将向量 \(f'(t)\) 记为 \(\mathrm{T}_t(f)(1) \in \mathrm{T}_{f(t)}(\mathbf{X})\)；该向量有时称为 \(f\) 在时刻 \(t\) 的速度。当 X 是 Banach 空间时，\(f'(t)\) 是 \(f\) 在 \(t\) 处的导数。

若有下式，则称 \(f\) 是 \(\xi\) 的积分曲线：

\[
f ^ {\prime} (t) = \xi (f (t))\tag{1}
\]

对所有 \(t \in I\) 都成立。这样的曲线属于 \(C^{s+1}\) 类。

若 \(x\) 是 X 的一点，则将满足 \(\xi\) 且 \(x\) 的 \(f: I \to X\) 的积分曲线 \(\xi\) 称为以 \(0 \in I\) 为起点的 \(f(0) = x\) 的积分曲线；对每个 \(x \in X\)，都存在一条以 \(\xi\) 为起点的 \(x\) 的积分曲线，它定义在 K 中 0 的某个适当开邻域上。

若 \(f: \mathbf{I} \to \mathbf{X}\) 是 \(\xi\) 的积分曲线且 \(a \in \mathbf{K}\)，则从 \(t \mapsto f(t - a)\) 到 \(\mathbf{I} + a\) 的映射 \(\mathbf{X}\) 是 \(\xi\) 的积分曲线。

设 \(f_{1}\colon\mathbf{I}_{1}\to\mathbf{X}\) 和 \(f_{2}\colon\mathbf{I}_{2}\to\mathbf{X}\) 是 \(\xi\) 的两条积分曲线，且 \(t\in\mathbf{I}_{1}\cap\mathbf{I}_{2}\)。若 \(f_{1}(t)=f_{2}(t)\)，则映射 \(f_{1}\) 和 \(f_{2}\) 在 \(t\) 的某个邻域内一致。

9.1.2（“流”）。设 W 是 X × K 的开子集。对每个 x ∈ X，记 W \(_{x}\) 为满足 (x, t) ∈ W 的 t ∈ K 所成的集合。称定义域为 W 的 ξ 的一个积分流为从 W 到 X 的 C \(^{k}\) 类映射 f（k ≤ r），满足：对每个 x ∈ X，由 f \(_{x}\) (t) = f(x, t) 定义的映射 f \(_{x}\) : W \(_{x}\) → X 是 ξ 的积分曲线，并且只要 (x, 0) ∈ W，就有 f(x, 0) = x。

对每个点 \(x \in X\)，存在一个 \(\xi\) 的 \(\mathbf{C}^s\) 类积分流，其定义域是 \((x, 0)\) 中 \(X \times K\) 的某个邻域。定义在 \(\xi\) 的某个邻域内的两个 \((x, 0)\) 的积分流在 \((x, 0)\) 的某个邻域内一致。

设 \(x\in X\)，且 \(f\) 是 \(\xi\) 的一个积分流，其定义域 W 是 \((x,0)\) 的一个邻域。存在 \((x,0,0)\) 在 \(X \times K \times K\) 中的一个邻域 V，使得若 \((x,t_{1},t_{2})\in\mathrm{V}\)，则有

\[
(x, t _ {1}) \in \mathbf {W}, \quad (f (x, t _ {1}), t _ {2}) \in \mathbf {W}, \quad (x, t _ {1} + t _ {2}) \in \mathbf {W}
\]

并且

\[
f (f (x, t _ {1}), t _ {2}) = f (x, t _ {1} + t _ {2}).\tag{2}
\]

若 X 是仿紧的，或者 X 是分离的且 K = R 或 C，则存在一个 \(\xi\) 的积分流，其定义域是 \(X \times \{0\}\) 的一个邻域。

9.1.3（“实情形”）。假设 K = R 且 X 是分离的。将定义域为 R 的开区间的 \(\xi\) 的任意积分曲线称为 \(\xi\) 的积分弧。若 \(f_{1}\) 的两条积分弧 \(f_{2}\) 和 \(\xi\) 分别定义在区间 \(I_{1}\) 和 \(I_{2}\) 上，并在 \(I_{1} \cap I_{2}\) 的一点处相同，则它们在 \(I_{1} \cap I_{2}\) 上一致。

对每个点 \(x \in X\)，存在一条以 \(x\) 为起点的积分弧，并且唯一存在 \(f_x: I_x \to X\)，使得对于每条积分弧 \(f: I \to X\)（以 \(x\) 为起点），都有 \(I \subset I_x\) 且 \(f = f_x|I\)。弧 \(f_x\) 称为向量场 \(\xi\) 以 \(x\) 为起点的最大积分弧，其定义域 \(I_x\) 有时称为点 \(x\) 对向量场 \(\xi\) 的寿命区间。

若 \(I_x = ]-a_-,a_+[\)（且若 \(a_+\) 有限，则 \(f_x\) 在区间 \([0,a_+[\) 上的限制是从 \([0,a_+[\) 到 X 的真映射），则特别地，\(f_x(t)\) 当 \(t\) 趋于 \(a_+\) 时没有极限。

9.1.4. 采用 9.1.3 的记号和假设，令 \(\Omega\) 为满足 \((x,t)\in\mathbf{X}\times\mathbf{R}\) 的对 \(t\in\mathbf{I}_{x}\) 所成的集合。它是 \(\mathbf{X}\times\mathbf{R}\) 中的开集，且由 \(f\colon\Omega\to\mathbf{X}\) 定义的映射 \(f(x,t)=f_{x}(t)\) 是定义域为 \(\xi\) 的 \(\Omega\) 的积分流。

由 \(\alpha_{+}\) 定义的函数 \(\alpha_{-}: X \to ]0,+\infty]\) 和 \(I_{x} = ]-\alpha_{-}(x),\alpha_{+}(x)[\) 是下半连续的。对于 \((x, t) \in \Omega\) 且 \(y = f(x, t)\)，有 \(\alpha_{+}(y) = \alpha_{+}(x) - t\) 和 \(\alpha_{-}(y) = \alpha_{-}(x) + t\)。

若 \(t_1 \in \mathbf{I}_x\) 且 \(t_2 \in \mathbf{I}_{f(x, t_1)}\)，则 \(t_1 + t_2 \in \mathbf{I}_x\)，并且

\[
f (f (x, t _ {1}), t _ {2}) = f (x, t _ {1} + t _ {2}).\tag{3}
\]

若函数 \(\alpha_{+}\)（分别为 \(\alpha_{-}\)）在 X 上有一个 \(>0\) 的常数作为下界，则它是常数且等于 \(+\infty\)。

在下列每种情形中，都有 \(\alpha_{+} = \alpha_{-} = +\infty\)（换言之，有 \(\Omega = \mathbf{X} \times \mathbf{R}\)）：

\begin{enumerate}
\def\labelenumi{(\roman{enumi})}
\item
  向量场 \(\xi\) 具有紧支集（例如 X 是紧的）；
\item
  存在一个保持 \(\xi\) 的 X 的传递自同构群；
\item
  X 上存在一个黎曼结构，使得 X 是完备的且 \(\xi\) 有界。
\end{enumerate}

若 \(\Omega = X \times R\)，则映射 \(f: X \times R \to X\) 是 \(C^s\) 类的群流形 \(R\) 在流形 \(X\) 上的右作用律，cf.~5.12.5。

9.1.5（“复情形”）。假设 K = C 且 X 是分离的。对于每个 \(a \in ]0, +\infty]\)，记 \(D_a\) 为 C 中以 0 为中心、半径为 a 的开圆盘。对于每个 \(x \in X\)，存在唯一的数 \(\rho(x) \in ]0, +\infty]\) 以及唯一的积分曲线 \(f_x: D_{\rho(x)} \to X\)，其起点为 \(x\)，使得对于每个 \(a \in ]0, +\infty]\) 以及每条积分曲线 \(f: \mathrm{D}_a \to \mathrm{X}\)，其起点为 \(x\)，都有 \(a \leqslant \rho(x)\) 且 \(f = f_x | \mathrm{D}_a\)。

对于每个 \(\lambda\in\mathbf{C}\)，令 \(\alpha_{\lambda}(x)\) 为点 \(x\) 对于实流形上的向量场 \(\lambda\xi\) 的寿命区间的上界。则有 \(\rho(x)=\inf_{|\lambda|=1}\alpha_{\lambda}(x)\)。

满足 \(\Delta\) 的对 \((x,t)\in X\times\mathbf{C}\) 所成的集合 \(t\in D_{\rho(x)}\) 是 \(X\times\mathbf{C}\) 中的开集，且由 \(f:\Delta\to X\) 定义的映射 \(f(x,t)=f_x(t)\) 是 \(\xi\) 的积分流。函数 \(\rho:X\to ]0,+\infty]\) 是下半连续的。对于 \((x,t)\in\Delta\) 且 \(y=f(x,t)\)，有 \(\rho(y)\geqslant\rho(x)-|t|\)。若 \(\rho\) 在 \(X\) 上有一个 \(>0\) 的常数作为下界，则 \(\rho = +\infty\) 且 \(\Delta = X \times C\)。

当 \(\Delta = X \times C\) 时，映射 \(f: X \times C \to X\) 是群流形 \(C\) 在流形 \(X\) 上的右作用律。

9.1.6. 设 \(\varphi: X \to Y\) 是流形的态射，\(\eta\) 是 Y 上的向量场，且 \(\xi\) 通过 \(\varphi\) 与 \(\eta\) 相关（参见 8.2.6）。若 \(f: I \to X\) 是 \(\xi\) 的积分曲线，则映射 \(\varphi \circ f: I \to Y\) 是 \(\eta\) 的积分曲线。若 \(K = R\) 且 X 和 Y 是分离的，则对于每个 \(x \in X\)，点 \(x\) 对向量场 \(\xi\) 的寿命区间包含于点 \(\varphi(x)\) 对向量场 \(\eta\) 的寿命区间。若此外 \(\varphi\) 是真映射，则这两个区间相等。

9.1.7（“含时方程”）。设 W 是 X × K 的开子集，且 Ξ: W → T(X) 是 pr \(^{s}\) : W → X 的 C \(_{1}\) 类提升，即从 W 到 T(X) 的 C \(^{s}\) 类映射，并满足对所有 (x, t) ∈ W，Ξ(x, t) 属于 T \(_{x}\) (X)。将从 K 的开子集 I 到 X 的 C \(^{k}\) 类映射 f 称为 Ξ 的积分曲线，其中 k ∈ N \(_{K}\) 且 k ≤ r，并且对所有 t ∈ I 都有

\[
(f (t), t) \in \mathrm{W},\qquad f ^ {\prime} (t) = \Xi (f (t), t).
\]

设 \(\eta\) 为 W 上的 \(C^s\) 类向量场，它由 \(\eta(x, t) = (\Xi(x, t), (t, 1))\) 定义。对于从 K 的开子集 I 到 X 的映射 \(f\)，要使其成为 \(\Xi\) 的积分曲线，充要条件是 \(t \mapsto (f(t), t)\) 按 9.1.1 的意义是 \(\eta\) 的积分曲线。

设 \(g: \mathbf{I} \to \mathbf{W}\) 是 \(\eta\) 的积分曲线，且 \(t_0\) 是 \(\mathbf{I}\) 中满足 \(\operatorname{pr}_2(g(t_0)) = t_0\) 的一点。满足 \(t \in \mathbf{I}\) 的 \(\operatorname{pr}_2(g(t)) = t\) 所成的集合是 \(t_0\) 的一个邻域；若 \(\mathbf{I}\) 连通，则该集合与 \(\mathbf{I}\) 重合。

9.1.8（“参数与初始条件”）。设 Z 是 \(C^{r}\) 类流形，W 是 X × K × Z 的开子集，且 \(\Xi\) 是 \(C^{s}\) 的 \(pr_{1}: W \to X\) 类提升。若 \(z \in Z\)，记满足 \(W_{z}\) 的 \((x, t) \in X \times K\) 所成的集合为 \((x, t, z) \in W\)，并将从 \(\Xi_{z}\) 到 T(X) 的映射 \((x, t) \mapsto \Xi(x, t, z)\) 记为 \(W_{z}\)。

设 \((x_0, t_0, z_0)\) 是 W 的一点。则 W 中存在 \(X_1 \times I_1 \times Z_1\) 的一个开邻域 \((x_0, t_0, z_0)\) 以及一个 \(C^s\) 类映射

\[
f \colon X _ {1} \times I _ {1} \times I _ {1} \times Z _ {1} \rightarrow X
\]

使得对于任意 \((x_{1}, t_{1}, z_{1}) \in \mathbf{X}_{1} \times \mathbf{I}_{1} \times \mathbf{Z}_{1}\)，映射 \(t \mapsto f(x_{1}, t_{1}, t, z_{1})\) 是 \(\Xi_{z_{1}}\) 的积分曲线（按 9.1.7 的意义），并满足关系 \(f(x_{1}, t_{1}, t_{1}, z_{1}) = x_{1}\)。

两个映射

\[
f _ {1}: \mathrm{X} _ {1} \times \mathrm{I} _ {1} \times \mathrm{I} _ {1} \times \mathrm{Z} _ {1} \rightarrow \mathrm{X} \quad \text {and} \quad f _ {2}: \mathrm{X} _ {2} \times \mathrm{I} _ {2} \times \mathrm{I} _ {2} \times \mathrm{Z} _ {2} \rightarrow \mathrm{X}
\]

满足上述条件的映射在 \((x_{0}, t_{0}, t_{0}, z_{0})\) 的某个邻域内一致。

